\section{Related work}
\label{sec:related}

\paragraph{Training-free high-resolution generation.}
Early work targeted convolutional U-Nets. ScaleCrafter~\citep{he2024scalecrafter} re-dilates convolutions, FouriScale~\citep{huang2024fouriscale} combines dilation with a low-pass constraint in the frequency domain, and HiDiffusion~\citep{zhang2024hidiffusion} reshapes U-Net features and windows its self-attention. DemoFusion~\citep{du2024demofusion}, FreeScale~\citep{qiu2024freescale}, and FreCaS~\citep{zhang2025frecas} generate progressively, from a low-resolution pass upward, and differ in how they fuse scales and schedule frequencies. For rectified-flow transformers, I-Max~\citep{du2024imax} projects the high-resolution flow toward a low-resolution guide and HiFlow~\citep{bu2025hiflow} aligns the high-resolution flow with a reference flow. Diffusion-4K~\citep{zhang2025diffusion4k} instead fine-tunes latent DiTs on its Aesthetic-4K set, which we use for evaluation prompts and references. Constancy RoPE is a single-pass change to attention and can in principle be combined with cascaded samplers.

\paragraph{RoPE extrapolation.}
RoPE~\citep{su2024roformer} encodes position as rotations at geometrically spaced frequencies. For language models, position interpolation~\citep{chen2023pi} divides all frequencies by the length ratio, NTK-aware scaling raises the base, YaRN~\citep{peng2024yarn} interpolates only the low-frequency dimensions and adds a global temperature, and LongRoPE~\citep{ding2024longrope} searches non-uniform per-dimension factors. These methods change frequencies, not magnitudes. In diffusion transformers, DyPE~\citep{issachar2025dype} writes the spectrum of $x_t$ per frequency and uses it to motivate a time-dependent RoPE schedule, and UltraImage~\citep{ultraimage2025} fixes the RoPE frequency whose period matches the training size and adds entropy-guided per-head focus. RIFLEx~\citep{zhao2025riflex} and UltraViCo~\citep{ultravico2025} study length extrapolation in video DiTs, where repetition appears in time. Barbero et al.~\citep{barbero2025round} show that high RoPE frequencies support positional attention patterns while low frequencies carry semantic content, which offers one explanation for why boosting local bands helps at high resolution.

\paragraph{Attention temperature and per-band gains.}
Jin et al.~\citep{jin2023training} keep attention entropy approximately constant across image sizes by scaling logits with the ratio of log token counts, and Scalable-Softmax~\citep{nakanishi2025ssmax} multiplies logits by a factor proportional to $\log n$. SEGA~\citep{sega2026} combines a power-law temperature $s^{0.16}$ on the logits with a per-RoPE-dimension modulation computed from latent spectra on top of a strengthened NTK base. AReA~\citep{area2026}, from the SEGA authors, moves from coarse cross-scale attention early in sampling to shrinking local windows later. To our knowledge, SEGA is the only prior training-free method that sets a per-band magnitude gain from a measured spectrum. Section~\ref{sec:sega} analyzes it in detail.

\paragraph{Spectral views of diffusion.}
Natural images have power spectra that fall roughly as a power of frequency~\citep{field1987relations}, and additive Gaussian noise is white, so diffusion fills in frequencies from coarse to fine~\citep{dieleman2024spectral,falck2025fourier}. Kadkhodaie et al.~\citep{kadkhodaie2024generalization} relate the generalization of diffusion denoisers to geometry-adaptive harmonic bases. Spectral Correction Guidance~\citep{scg2026} applies a per-step, per-radial-bin gain to latent magnitudes toward a fitted spectral law at $1024^2$; it acts on the latent, not on attention. CSFlow~\citep{csflow2026} reweights diffusion timesteps so that the frequencies being generated follow the human contrast sensitivity function. Constancy RoPE shares the per-band signal-to-noise view with these works and is, to our knowledge, the first to derive a per-band attention gain from it and to close the loop on the model's own native-resolution spectrum.

\paragraph{Human spatial vision.}
Contrast sensitivity is band-pass in daylight and low-pass in dim light~\citep{vannes1967spatial}, which efficient-coding theory explains as whitening combined with Wiener filtering under noise~\citep{atick1990towards,atick1992retina,vanhateren1992theory}. Above threshold, perceived contrast is approximately constant across spatial frequency, so the visual system compensates for the blur of its own optics~\citep{georgeson1975contrast}. Adaptation to blurred or sharpened images shifts perceived focus toward natural spectral slopes~\citep{webster2002neural}, and observers discriminate spectral slope best near natural values~\citep{knill1990human,parraga2000human}. Spatial vision analyzes images in channels of roughly constant octave bandwidth~\citep{campbell1968application,field1987relations}, and cortical responses are divided by the pooled activity of their neighbors~\citep{heeger1992normalization,carandini2012normalization,schwartz2001natural}. Natural video follows a joint space-time spectral law~\citep{dong1995statistics}, which suggests an extension to temporal RoPE bands that we do not test here.
